\chapter{Chorea}

\section*{Definition}
Chorea is a hyperkinetic movement disorder characterized by abrupt, irregular, unpredictable, non-rhythmic, flowing movements that migrate randomly from one body part to another. Movements are involuntary but may be partially incorporated into voluntary actions (parakinesia), distinguishing chorea from other involuntary movements.

\section*{What Happens in Your Body}
Chorea results from dysfunction within the basal ganglia, particularly involving reduced GABAergic output from the striatum to the external globus pallidus, leading to excessive thalamocortical drive. Imbalance between the direct and indirect basal ganglia pathways, especially decreased activity of the indirect pathway via striatal enkephalinergic neurons, disinhibits thalamic excitation of the motor cortex. Altered dopamine signaling, glutamatergic excitotoxicity, and mitochondrial dysfunction contribute to neuronal vulnerability.

\section*{Causes}
\begin{itemize}
  \item Genetic causes: Huntington disease (CAG repeat expansion in HTT), benign hereditary chorea (NKX2-1), neuroacanthocytosis, dentatorubral-pallidoluysian atrophy (DRPLA), Wilson disease, ataxia-telangiectasia, Friedreich ataxia, pantothenate kinase-associated neurodegeneration
  \item Vascular causes: Ischemic or hemorrhagic stroke involving caudate, putamen, or thalamus; cerebral vasculitis
  \item Infectious and postinfectious: Sydenham chorea (group A streptococcal infection, rheumatic fever), HIV, toxoplasmosis, Lyme disease, Creutzfeldt-Jakob disease
  \item Metabolic and endocrine: Hyperthyroidism, hyperglycemia (hemichorea-hemiballism in nonketotic hyperglycemia), hypocalcemia, hypernatremia, hepatic encephalopathy, pregnancy (chorea gravidarum)
  \item Drug-induced: Levodopa (wearing-off dyskinesias in Parkinson disease), antipsychotics (tardive dyskinesia), oral contraceptives, anticonvulsants (phenytoin), stimulants, alcohol withdrawal
  \item Autoimmune: Anti-phospholipid syndrome, systemic lupus erythematosus, anti-NMDA receptor encephalitis, paraneoplastic syndromes (anti-CRMP5, anti-Hu)
\end{itemize}

\section*{When to See a Doctor}
See your doctor if:
\begin{itemize}
  \item New-onset chorea, especially with family history of movement disorders
  \item Progressive chorea with cognitive decline, psychiatric changes, or behavioral disturbance
  \item Acute onset of unilateral chorea (hemichorea) suggesting stroke or hyperglycemia
  \item Chorea in child with preceding sore throat or strep infection (Sydenham chorea)
  \item Associated signs of systemic disease (rash, arthritis, fever, weight loss, thyroid enlargement)
\end{itemize}

\section*{Related Symptoms}
\begin{itemize}
  \item Athetosis
  \item Ballism (hemiballismus)
  \item Dystonia
  \item Tics
  \item Myoclonus
  \item Gait instability
  \item Cognitive impairment
  \item Psychiatric disturbances
\end{itemize}
\textbf{Important:} Huntington disease typically presents with a triad of chorea, cognitive decline, and psychiatric disturbance, with age of onset inversely correlated with CAG repeat length; the UHDRS (Unified Huntington's Disease Rating Scale) quantifies disease severity.

\section*{Self-Care / Home Management}
\begin{itemize}
  \item Fall prevention through home safety modifications and use of mobility aids if gait is affected.
  \item Stress reduction as emotional stress and anxiety exacerbate choreiform movements.
  \item Nutritional support with calorie-dense foods, as hyperkinetic movements increase energy expenditure.
  \item Occupational therapy for adaptive strategies to compensate for involuntary movements during daily activities.
\end{itemize}

\section*{How Your Doctor Will Evaluate You}
\begin{itemize}
  \item History includes age of onset, family history, medication and drug exposure, associated neurologic and systemic symptoms.
  \item Neurologic examination assesses distribution, severity, and phenomenology of chorea along with cognitive and psychiatric evaluation.
  \item Genetic testing for Huntington disease (HTT CAG repeats) and other genetic choreas guided by phenotype.
  \item Laboratory evaluation includes thyroid function, glucose, electrolytes, liver function, rheumatologic panel (ANA, anti-phospholipid antibodies), and infectious serologies.
  \item Brain MRI evaluates for caudate atrophy (Huntington disease), strokes, or structural basal ganglia lesions.
\end{itemize}

\section*{Treatment Options}
\begin{itemize}
  \item Dopamine-depleting agents: Tetrabenazine, deutetrabenazine, and valbenazine are FDA-approved for Huntington chorea.
  \item Antipsychotics: Risperidone, olanzapine, haloperidol reduce chorea through dopamine receptor blockade.
  \item Benzodiazepines (clonazepam) and anticonvulsants (levetiracetam, topiramate) provide adjunctive benefit.
  \item Sydenham chorea is treated with antibiotics for streptococcal infection plus symptomatic therapy (valproate, carbamazepine, or neuroleptics) and secondary prophylaxis to prevent rheumatic recurrence.
  \item Underlying metabolic, endocrine, or autoimmune causes are treated specifically (thyroid hormone normalization, glucose control, immunosuppression).
\end{itemize}

\section*{References}
\begin{thebibliography}{99}
\bibitem{ref1} Pringsheim T, Okun MS, Muller-Vahl K, et al. "Practice Guideline: Treatment of Tics in People with Tourette Syndrome." Neurology. 2019;92(19):896--906.
\bibitem{ref2} Cardoso F, Seppi K, Mair KJ, et al. "Semiology of Chorea." Mov Disord Clin Pract. 2016;3(2):118--125.
\bibitem{ref3} Shannon KM. "Huntington Disease." In: Jankovic J, Tolosa E, eds. Parkinson's Disease and Movement Disorders. 6th ed. Lippincott Williams \& Wilkins; 2015.
\bibitem{ref4} Walker RH. "Differential Diagnosis of Chorea." Curr Neurol Neurosci Rep. 2011;11(4):385--395.
\bibitem{ref5} Bhidayasiri R, Truong DD. "Chorea and Related Disorders." Postgrad Med J. 2004;80(947):527--534.
\bibitem{ref6} Ross CA, Tabrizi SJ. "Huntington's Disease: From Molecular Pathogenesis to Clinical Treatment." Lancet Neurol. 2011;10(1):83--98.
\end{thebibliography}
